\section{Estrategias de estabilización y regularización}

\subsection{Bootstrapping y normalización de gradientes}

El bootstrapping vincula la salida de la red con su propio objetivo, lo que fácilmente provoca explosión de gradientes. En la práctica se suelen usar el recorte de gradientes y la normalización para controlarla.

\begin{importantbox}{Referencia rápida de técnicas de estabilización}
\begin{itemize}
    \item Recorte de gradientes a 1.0, para evitar que actualizaciones demasiado grandes provoquen divergencia.
    \item Layer normalization y weight decay controlan el rango de la salida.
    \item La separación entre selección de acción y estimación de valor de Double Q también es parte de las estrategias de estabilización.
\end{itemize}
\end{importantbox}

\subsection{Moldeo de recompensa y normas de comportamiento}

En tareas de alta dimensión, usar directamente el reward produce alta varianza. Se puede añadir un shaping term o una entropy penalty para suavizar la learning signal.

\begin{warningbox}{Las trampas del moldeo de recompensa}
El shaping term debe ser potential-based; de lo contrario puede cambiar la política óptima. Añadir una penalty tampoco debe permitir que el valor de Q crezca sin límite.
\end{warningbox}

\subsection{Restricción de acciones y recorte de salida}

Para evitar que la salida de la función Q se desvíe, se pueden usar clipping y output constraint:

\begin{itemize}
    \item Limitar el TD target al observed reward range ± margin, para evitar que la estimación de valor explote.
    \item Introducir Q regularization (como L2 shrinkage toward previous target) para que las actualizaciones sean más suaves.
    \item Aplicar clamp tras añadir ruido al muestreo de la action, para garantizar que $\arg\max_a Q(s,a)$ no caiga en una región no observada.
\end{itemize}

\subsection{Resumen de la sección}

La estabilización debe abordarse simultáneamente desde la señal del gradiente y la señal del comportamiento, para evitar tanto la explosión de la función de valor como la deriva de la política.

